\documentclass[11pt, letterpaper]{article}
\usepackage{mobi-lectura}

\title{Análisis de Salida: Simulaciones de Estado Estacionario}
\author{Alejandra Tabares Pozos}
\date{}

\begin{document}

\maketitle
\tableofcontents
\newpage

% ══════════════════════════════════════════════════
\section{¿Cuándo interesa el estado estacionario?}
% ══════════════════════════════════════════════════

En muchos sistemas, la pregunta de interés no es ``¿cuánto esperará un cliente durante el turno de hoy?'' sino ``¿cuánto espera un cliente típico en el largo plazo, bajo condiciones normales de operación?''. Este es el régimen de \textbf{estado estacionario}: el comportamiento del sistema cuando el efecto de las condiciones iniciales se ha disipado.

\begin{definicionbox}[Estado estacionario de una simulación]
Una simulación se dice que está en estado estacionario cuando la distribución de su proceso de estado $X(t)$ ya no depende del tiempo: $\Prob(X(t) \in A) \approx \pi(A)$ para todo conjunto medible $A$, donde $\pi$ es la distribución estacionaria.
\end{definicionbox}

\textbf{Ejemplos de uso estacionario}:
\begin{itemize}
    \item Fábrica con producción continua 24/7.
    \item Call center con demanda promedio constante.
    \item Red de distribución con demanda diaria estable.
    \item Comparación de alternativas de diseño del sistema.
\end{itemize}

El parámetro de interés es típicamente $\theta = \E_\pi[g(X)]$, el valor esperado de una función del proceso en régimen estacionario.

% ══════════════════════════════════════════════════
\section{El problema del período transitorio}
% ══════════════════════════════════════════════════

Toda simulación arranca desde un estado inicial $X(0)$. En la práctica, el estado inicial es conveniente (p.ej.\ sistema vacío), pero no necesariamente representativo del estado estacionario. Durante el \textbf{período transitorio} (o \emph{warm-up}), el proceso aún no ha alcanzado su distribución estacionaria.

Si se incluyen las observaciones del período transitorio en el cálculo de la media, se introduce un \textbf{sesgo de inicialización}:

\[
\E\!\left[\frac{1}{m}\sum_{i=1}^m Y_i\right] = \theta + \underbrace{\frac{1}{m}\sum_{i=1}^m\bigl(\E[Y_i] - \theta\bigr)}_{\text{sesgo de inicialización}}
\]

El sesgo depende de la distancia entre el estado inicial y la distribución estacionaria, y decrece con $m$, pero lentamente. Crucialmente, este sesgo \textbf{no desaparece al aumentar el número de réplicas}; solo desaparece si se descarta o controla el período transitorio.

\begin{center}
\begin{tikzpicture}
\begin{axis}[
  width=12cm, height=5cm,
  xlabel={Tiempo $t$ (unidades de simulación)},
  ylabel={$\bar Y(t)$ (media acumulada)},
  xmin=0, xmax=500, ymin=0, ymax=10,
  grid=both, grid style={line width=0.3pt, draw=gray!30},
  tick label style={font=\scriptsize},
  label style={font=\small},
  legend style={font=\scriptsize, at={(0.98,0.95)}, anchor=north east}
]
% Transient period decaying from 0 toward theta~6
\addplot[thick, mobiblue, smooth] coordinates {
  (0,0)(20,1.5)(40,3.0)(60,4.2)(80,5.0)(100,5.5)(120,5.8)(150,5.9)(200,6.05)
  (250,6.1)(300,6.05)(350,5.95)(400,6.1)(450,6.0)(500,6.05)
};
\addplot[thick, red!70!black, dashed] coordinates {(0,6)(500,6)};
\addplot[thick, mobigold!80!black, dotted] coordinates {(120,0)(120,10)};
\addlegendentry{$\bar Y(t)$ — media acumulada}
\addlegendentry{$\theta$ — valor estacionario real}
\addlegendentry{Fin de warm-up $t_w$}
\end{axis}
\end{tikzpicture}
\end{center}

La media acumulada parte de un valor bajo (sistema vacío) y converge hacia $\theta$. El período $[0, t_w]$ debe descartarse.

\begin{alertabox}[El sesgo no se promedia con réplicas]
Si se ejecutan $n$ réplicas \emph{sin} descartar el warm-up, todas empiezan desde el mismo estado inicial sesgado. El promedio entre réplicas reduce la varianza pero \emph{no} el sesgo: $\E[\bar Y] = \theta + \text{sesgo}$ para cualquier $n$. El sesgo se elimina descartando el transitorio, no aumentando réplicas.
\end{alertabox}

% ══════════════════════════════════════════════════
\section{Método de Welch para estimar $t_w$}
% ══════════════════════════════════════════════════

El método de Welch (1983) ofrece un criterio gráfico para identificar cuándo termina el período transitorio.

\begin{metodobox}[Método de Welch]
\begin{enumerate}
    \item Ejecutar $k \ge 5$ réplicas largas, todas desde el mismo estado inicial. Registrar la secuencia de observaciones $Y_{ij}$ (observación $i$ de la réplica $j$).
    \item Para cada $i$, calcular el promedio entre réplicas: $\bar Y_i = k^{-1}\sum_{j=1}^k Y_{ij}$.
    \item Suavizar con media móvil de ventana $w$: $\tilde Y_i = (2w+1)^{-1}\sum_{l=-w}^{w}\bar Y_{i+l}$.
    \item Graficar $\tilde Y_i$ vs.\ $i$ para distintos valores de $w$.
    \item Elegir $t_w$ como el punto a partir del cual $\tilde Y_i$ se estabiliza visualmente.
\end{enumerate}
\end{metodobox}

\textbf{Elección de la ventana $w$}:

\begin{itemize}
    \item Una ventana $w$ muy pequeña no elimina el ruido; la curva suavizada es ruidosa y difícil de interpretar.
    \item Una ventana $w$ muy grande oculta el punto de estabilización, suavizando demasiado.
    \item Regla práctica: probar $w \in \{n/10,\; n/20,\; n/40\}$ donde $n$ es la longitud de la serie. Elegir el mínimo $w$ con el que la curva sea claramente estable.
\end{itemize}

\begin{notabox}[Limitaciones del método de Welch]
El método de Welch es \textbf{subjetivo}: depende de la inspección visual del analista. Distintos analistas pueden elegir $t_w$ diferentes. Alternativas más formales incluyen:
\begin{itemize}
    \item \textbf{MSER-5} (Marginal Standard Error Rule): elige $t_w$ que minimiza el error estándar marginal de $\bar Y$ calculado sobre $\{Y_{t_w+1}, \ldots, Y_n\}$.
    \item \textbf{Test de Schruben}: prueba estadística que detecta si hay tendencia en la media.
\end{itemize}
En la práctica, Welch sigue siendo el método más usado por su simplicidad y transparencia.
\end{notabox}

% ══════════════════════════════════════════════════
\section{Método de Batch Means (medias de lote)}
% ══════════════════════════════════════════════════

El método de Batch Means es la técnica más común para estimación en estado estacionario. La idea es convertir una larga corrida autocorrelada en un conjunto de lotes aproximadamente independientes.

\subsection{Procedimiento}

\begin{metodobox}[Batch Means]
\begin{enumerate}
    \item Hacer un warm-up de longitud $t_w$ (descartar observaciones del transitorio).
    \item Ejecutar la simulación durante $n = m \cdot b$ observaciones adicionales.
    \item Dividir en $m$ lotes de $b$ observaciones cada uno:
    \[
    \bar Y_j = \frac{1}{b}\sum_{i=(j-1)b+1}^{jb} Y_i, \qquad j = 1,\ldots,m
    \]
    \item Tratar $\bar Y_1,\ldots,\bar Y_m$ como i.i.d.\ (válido si $b$ es suficientemente grande para que la autocorrelación entre lotes sea despreciable).
    \item Calcular IC estándar: $\bar Y_m \pm t_{\alpha/2,m-1}\cdot S_B/\sqrt{m}$, donde $S_B^2 = \frac{1}{m-1}\sum(\bar Y_j - \bar Y_m)^2$.
\end{enumerate}
\end{metodobox}

\subsection{¿Por qué funciona Batch Means?}

Sea $\sigma_{\bar Y}^2 = \Var[\bar Y_n]$ la varianza real de la media muestral (que incluye los efectos de autocorrelación). Para un proceso estacionario con autocovarianza $\gamma(k) = \text{Cov}(Y_i, Y_{i+k})$:

\[
\sigma_{\bar Y}^2 = \frac{1}{n}\left[\gamma(0) + 2\sum_{k=1}^{n-1}\left(1-\frac{k}{n}\right)\gamma(k)\right] \approx \frac{\sigma_0^2}{n}
\]

donde $\sigma_0^2 = \gamma(0) + 2\sum_{k=1}^{\infty}\gamma(k)$ se llama la \textbf{varianza de largo plazo}. El estimador de Batch Means $S_B^2/m$ converge a $\sigma_0^2/n$ cuando $b \to \infty$, es decir, estima correctamente la varianza de $\bar Y_n$ incluyendo la autocorrelación.

\subsection{Elección de $b$ y $m$}

\begin{itemize}
    \item $b$ debe ser grande en comparación con el tiempo de correlación de la serie. Una regla práctica: $b \ge 30\tau$ donde $\tau$ es el rezago donde el autocorrelograma decae por debajo de $0{,}05$.
    \item $m$ determina los grados de libertad del IC; se recomienda $m \ge 20$ (idealmente $m \ge 30$).
    \item Si $n = m \cdot b$ resulta demasiado largo, se puede reducir $m$ a 10--15 aceptando IC más anchos.
\end{itemize}

\subsection{Verificación de independencia de lotes}

\textbf{Razón de Von Neumann}. Se calcula:
\[
VN = \frac{\sum_{j=1}^{m-1}(\bar Y_{j+1} - \bar Y_j)^2}{(m-1)\,S_B^2}
\]

Bajo independencia, $\E[VN] \approx 2$. Si $VN \ll 2$, los lotes consecutivos están positivamente correlacionados: $b$ es insuficiente. En ese caso, aumentar $b$ (o equivalentemente, reducir $m$ manteniendo $n$).

\begin{alertabox}[Error frecuente: lotes demasiado cortos]
Si $b$ es insuficiente, $\bar Y_j$ y $\bar Y_{j+1}$ siguen estando correlacionados, lo que hace que $S_B^2$ subestime la varianza real y el IC sea demasiado estrecho (falsa precisión). Verificar siempre con la razón de Von Neumann o inspeccionando el autocorrelograma de $\{\bar Y_j\}$.
\end{alertabox}

\subsection{Batch Means con solapamiento (OBM)}

El método clásico (non-overlapping batch means, NBM) desperdicia información al dividir en lotes disjuntos. Los \textbf{overlapping batch means} (OBM) usan lotes solapados:

\[
\bar Y_j^{\text{OBM}} = \frac{1}{b}\sum_{i=j}^{j+b-1} Y_i, \qquad j = 1,\ldots,n-b+1
\]

OBM produce un estimador de varianza con menor error cuadrático medio que NBM para el mismo $n$ y $b$, a costa de mayor complejidad computacional. En la práctica, OBM reduce la varianza del estimador de varianza en un factor de aproximadamente $2/3$ respecto a NBM.

% ══════════════════════════════════════════════════
\section{Método de réplicas con calentamiento (\emph{deletion method})}
% ══════════════════════════════════════════════════

Alternativa al Batch Means: ejecutar $n$ réplicas independientes, descartar las primeras $d$ observaciones de cada réplica (warm-up) y promediar el resto.

\[
Y_j = \frac{1}{m_j - d}\sum_{i=d+1}^{m_j} Y_{ij}, \qquad j = 1,\ldots,n
\]

Las $Y_j$ son i.i.d.\ y se aplica la misma inferencia $t$ del Tema 3.3. La ventaja es la simplicidad; la desventaja es el desperdicio de las $n \cdot d$ observaciones del warm-up (una por réplica).

\begin{notabox}[¿Batch Means o réplicas con calentamiento?]
\begin{itemize}
    \item \textbf{Batch Means}: más eficiente computacionalmente (una sola corrida larga, un solo warm-up). Requiere verificar independencia de lotes.
    \item \textbf{Réplicas con calentamiento}: más simples de implementar y verificar. Múltiples warm-ups desperdician cómputo. Facilitan la sincronización de CRN (Sección 7).
\end{itemize}
En general, Batch Means se prefiere para estimación de una sola métrica; réplicas con calentamiento se prefieren para comparaciones con CRN.
\end{notabox}

% ══════════════════════════════════════════════════
\section{Comparación de dos alternativas}
% ══════════════════════════════════════════════════

Sea $\theta^{(1)}$ y $\theta^{(2)}$ la métrica de interés para dos configuraciones del sistema (p.ej.\ 2 vs.\ 3 servidores). Se estima la diferencia $\Delta = \theta^{(1)} - \theta^{(2)}$.

\subsection{Muestras independientes}

Si se ejecutan $n_1$ réplicas del sistema 1 y $n_2$ réplicas del sistema 2 con semillas independientes:
\[
\hat\Delta = \bar Y_{n_1}^{(1)} - \bar Y_{n_2}^{(2)}, \qquad
\Var[\hat\Delta] = \frac{\sigma_1^2}{n_1} + \frac{\sigma_2^2}{n_2}
\]

El IC de $\Delta$ es:
\[
\hat\Delta \;\pm\; t_{\alpha/2,\,\nu} \cdot \sqrt{\frac{S_1^2}{n_1} + \frac{S_2^2}{n_2}}
\]

donde $\nu$ se calcula con la fórmula de Welch-Satterthwaite. Si el IC no contiene 0, hay diferencia significativa.

\subsection{Números aleatorios comunes (CRN)}

Los \textbf{CRN} (Common Random Numbers) son una técnica de reducción de varianza: se usa la \emph{misma} secuencia de números aleatorios en ambas corridas (mismas llegadas, misma asignación de tiempos de servicio). Esto introduce correlación \emph{positiva} entre $Y_j^{(1)}$ y $Y_j^{(2)}$.

\textbf{¿Por qué CRN reduce la varianza?} La varianza de la diferencia pareada es:

\[
\Var[\Delta_j] = \Var[Y_j^{(1)}] + \Var[Y_j^{(2)}] - 2\,\text{Cov}(Y_j^{(1)},Y_j^{(2)})
\]

Con CRN, los sistemas comparten la misma aleatoriedad de entrada, por lo que $\text{Cov}(Y_j^{(1)},Y_j^{(2)}) > 0$. Esto reduce $\Var[\Delta_j]$ respecto al caso independiente donde $\text{Cov} = 0$. La reducción es mayor cuando los dos sistemas son similares (lo cual hace que la correlación sea alta).

\textbf{Condición necesaria}: CRN solo funciona si $\text{Corr}(Y_j^{(1)}, Y_j^{(2)}) > 0$, lo cual requiere \textbf{sincronización correcta} de los números aleatorios.

\begin{metodobox}[Implementación de CRN con réplicas pareadas]
\begin{enumerate}
    \item Asignar un flujo de números aleatorios distinto a cada proceso estocástico: un flujo para llegadas, otro para servicios tipo A, otro para servicios tipo B, etc.
    \item En cada réplica $j$, usar los \emph{mismos} flujos para la simulación 1 y la simulación 2.
    \item Calcular $\Delta_j = Y_j^{(1)} - Y_j^{(2)}$ para cada réplica.
    \item Aplicar inferencia $t$ sobre la muestra $\{\Delta_j\}$ (prueba de una muestra sobre $H_0: \Delta = 0$):
    \[
    IC_{1-\alpha}(\Delta) = \bar\Delta_n \pm t_{\alpha/2,n-1} \cdot \frac{S_\Delta}{\sqrt{n}}
    \]
\end{enumerate}
\end{metodobox}

\begin{alertabox}[CRN puede aumentar la varianza si se usa incorrectamente]
Si los números aleatorios de los dos sistemas no están correctamente sincronizados (mismo número aleatorio para procesos distintos), CRN puede producir correlación negativa y \emph{aumentar} la varianza en lugar de reducirla. La sincronización requiere que el mismo número aleatorio se use para el mismo ``tipo'' de evento en ambos sistemas. Usar flujos separados (streams) es la técnica estándar.
\end{alertabox}

% ══════════════════════════════════════════════════
\section{Comparación de $k > 2$ alternativas}
% ══════════════════════════════════════════════════

Cuando se comparan más de dos configuraciones, el problema se complica: hacer $\binom{k}{2}$ comparaciones por pares con nivel $\alpha$ individual da un nivel global mucho mayor que $\alpha$ (problema de comparaciones múltiples).

\subsection{Corrección de Bonferroni}

Para mantener un nivel de confianza global $1-\alpha$ con $k$ sistemas (y $\binom{k}{2}$ comparaciones), usar nivel individual $\alpha' = \alpha/\binom{k}{2}$ en cada comparación. Esto es conservador pero simple.

\subsection{Ranking and Selection}

Los procedimientos de \textbf{ranking and selection} (R\&S) están diseñados específicamente para seleccionar el mejor sistema (el de menor $W_q$, por ejemplo) con una garantía probabilística.

\begin{definicionbox}[Procedimiento de selección]
Dado un conjunto de $k$ sistemas y un parámetro de indiferencia $\delta^*$ (la menor diferencia que el analista considera prácticamente significativa), un procedimiento R\&S garantiza:
\[
\Prob(\text{seleccionar el mejor sistema}) \ge 1-\alpha
\]
siempre que el mejor sistema sea al menos $\delta^*$ mejor que el segundo mejor.
\end{definicionbox}

El \textbf{procedimiento KN} (Kim \& Nelson, 2006) es el más usado: es secuencial (añade réplicas según se necesiten) y permite eliminación temprana de alternativas claramente inferiores.

\begin{notabox}[Cuándo usar cada enfoque]
\begin{itemize}
    \item \textbf{2 sistemas}: prueba $t$ pareada con CRN (Sección 6.2).
    \item \textbf{3--5 sistemas}: comparaciones por pares con Bonferroni y CRN.
    \item \textbf{$>5$ sistemas}: procedimientos de R\&S (KN) para eficiencia.
\end{itemize}
\end{notabox}

% ══════════════════════════════════════════════════
\section{Ejemplo básico: cola $M/M/1$ en estado estacionario}
% ══════════════════════════════════════════════════

\begin{ejemplobox}[Sistema]
Se modela una línea de producción como cola $M/M/1$ con $\lambda=4$ piezas/min y $\mu=5$ piezas/min ($\rho=0{,}8$). El sistema opera 24/7 y se quiere estimar el tiempo de espera en cola en régimen estacionario. Se usa una sola corrida larga con Batch Means.
\end{ejemplobox}

\textbf{Referencia analítica}:
\[
W_q = \frac{\rho}{\mu(1-\rho)} = \frac{0{,}8}{5 \times 0{,}2} = 0{,}80 \text{ min}
\]

\subsection{Determinación del warm-up (Welch)}

Se ejecutan $k=10$ réplicas de 10\,000 clientes cada una. Se calcula $\bar Y_i$ (promedio del tiempo de espera del cliente $i$-ésimo entre las 10 réplicas) y se suaviza con ventana $w=200$.

La curva suavizada parte de $\approx 0$ (primer cliente nunca espera en sistema vacío), crece hasta estabilizarse en $\approx 0{,}8$ min alrededor del cliente 3\,000. Se elige conservadoramente $t_w = 5\,000$ clientes.

\subsection{Batch Means con una corrida larga}

Se ejecuta una \textbf{sola corrida} de 50\,000 clientes. Se descartan los primeros 5\,000 (warm-up). Los 45\,000 restantes se dividen en $m=30$ lotes de $b=1\,500$ clientes cada uno.

\[
\bar Y_j = \frac{1}{1500}\sum_{i=(j-1)\cdot 1500+1}^{j\cdot 1500} W_i, \qquad j=1,\ldots,30
\]

\textbf{Verificación de independencia}: la razón de Von Neumann da $VN = 1{,}89$. Como $1{,}89 \approx 2$, los lotes son aproximadamente independientes. $\checkmark$

\textbf{Resultados}:
\[
\bar Y_{30} = 0{,}814 \text{ min}, \quad S_B = 0{,}098 \text{ min}
\]
\[
IC_{95\%} = 0{,}814 \pm 2{,}045 \times \frac{0{,}098}{\sqrt{30}} = 0{,}814 \pm 0{,}037 = [0{,}777;\; 0{,}851] \text{ min}
\]

El IC contiene el valor analítico $0{,}80$. Error relativo: $|0{,}814 - 0{,}80|/0{,}80 = 1{,}8\%$.

\begin{resultadobox}[Resultado del ejemplo básico]
Usando Batch Means con una corrida de 50\,000 clientes (5\,000 de warm-up, 30 lotes de 1\,500), el tiempo de espera estacionario se estima en $\bar W_q = 0{,}814$ min con IC 95\% $[0{,}777;\; 0{,}851]$ min. La razón de Von Neumann confirma la independencia de lotes, y el IC contiene la referencia analítica.
\end{resultadobox}

% ══════════════════════════════════════════════════
\section{Ejemplo intermedio: comparar $M/M/1$ vs.\ $M/D/1$ con CRN}
% ══════════════════════════════════════════════════

\begin{ejemplobox}[Sistema]
Continuando con el sistema del ejemplo básico ($\lambda=4$, $\mu=5$, $\rho=0{,}8$), se evalúan dos líneas de producción con la misma carga pero diferente variabilidad de servicio:
\begin{itemize}
    \item \textbf{Sistema A ($M/M/1$)}: tiempos de servicio exponenciales, $\text{CV}_S=1$.
    \item \textbf{Sistema B ($M/D/1$)}: servicio determinista (robot), $S=0{,}2$ min exactos, $\text{CV}_S=0$.
\end{itemize}
Se quiere estimar la diferencia en $W_q$ y determinar si es significativa, usando CRN para máxima eficiencia.
\end{ejemplobox}

\textbf{Referencias analíticas} (P-K):
\[
W_q^A = \frac{\rho}{\mu(1-\rho)} = 0{,}80 \text{ min}
\]
\[
W_q^B = \frac{\lambda\E[S^2]}{2(1-\rho)} = \frac{4 \times 0{,}04}{2 \times 0{,}2} = 0{,}40 \text{ min}
\]
El servicio determinista reduce el tiempo en cola exactamente a la mitad (factor $\frac{1+\text{CV}_S^2}{2} = 0{,}5$).

\subsection{Diseño del estudio}

Se usan réplicas con calentamiento (deletion method): $n=30$ réplicas de 6\,000 clientes cada una, descartando los primeros 1\,000 (warm-up determinado por Welch en el ejemplo básico). CRN: las mismas llegadas alimentan ambos sistemas en cada réplica.

\subsection{Resultados individuales}

\begin{center}\small
\renewcommand{\arraystretch}{1.3}
\begin{tabular}{lccc}
\toprule
\textbf{Sistema} & $\bar W_q$ (min) & IC 95\% & Ref.\ analítica \\
\midrule
A ($M/M/1$) & $0{,}812$ & $[0{,}783;\;0{,}841]$ & $0{,}800$ $\checkmark$ \\
B ($M/D/1$) & $0{,}404$ & $[0{,}391;\;0{,}417]$ & $0{,}400$ $\checkmark$ \\
\bottomrule
\end{tabular}
\end{center}

Ambos IC contienen la referencia analítica.

\subsection{Diferencia pareada con CRN}

Para cada réplica $j$: $\Delta_j = W_{q,j}^{A} - W_{q,j}^{B}$.

\[
\bar\Delta = 0{,}408 \text{ min}, \quad S_\Delta = 0{,}061 \text{ min}
\]
\[
IC_{95\%}(\Delta) = 0{,}408 \pm 2{,}045 \times \frac{0{,}061}{\sqrt{30}} = 0{,}408 \pm 0{,}023 = [0{,}385;\; 0{,}431] \text{ min}
\]

El IC no contiene 0: la diferencia es \textbf{estadísticamente significativa}. El servicio determinista reduce la espera en $\approx 0{,}41$ min.

\subsection{Eficiencia de CRN}

\begin{center}\small
\renewcommand{\arraystretch}{1.25}
\begin{tabular}{lcc}
\toprule
\textbf{Método} & $S_\Delta$ (min) & Semiancho IC (min) \\
\midrule
Muestras independientes & $0{,}198$ & $0{,}074$ \\
CRN (réplicas pareadas) & $0{,}061$ & $0{,}023$ \\
\bottomrule
\end{tabular}
\end{center}

CRN redujo el semiancho del IC en un \textbf{69\%}. Equivalentemente, para lograr el mismo IC con muestras independientes se necesitarían:
\[
n_{\text{indep}} = \left\lceil 30 \times \left(\frac{0{,}198}{0{,}061}\right)^2\right\rceil = \left\lceil 30 \times 10{,}5\right\rceil = 316 \text{ réplicas}
\]

CRN ahorró un factor de $\approx 10\times$ en cómputo.

\begin{resultadobox}[Resultado del ejemplo intermedio]
La comparación pareada con CRN demuestra que el servicio determinista reduce $W_q$ en $0{,}41$ min (IC $[0{,}385;\; 0{,}431]$). CRN fue altamente eficiente: redujo la varianza del estimador de la diferencia por un factor de 10. Esto es coherente con que los dos sistemas comparten toda la aleatoriedad de llegadas y difieren solo en la distribución de servicio.
\end{resultadobox}

% ══════════════════════════════════════════════════
\section{Ejemplo escalado: comparación de políticas en un call center}
% ══════════════════════════════════════════════════

\begin{ejemplobox}[Sistema]
Un call center opera 24/7 con demanda promedio constante. Las llamadas llegan con tasa $\lambda=45$ llamadas/h (proceso de Poisson). El tiempo de atención sigue una distribución \textbf{Log-normal}($\mu_L=1{,}7$, $\sigma_L=0{,}6$), lo que da:
\[
\E[S] = e^{1{,}7+0{,}18} = e^{1{,}88} \approx 6{,}55 \text{ min}, \quad
\text{CV}_S = \sqrt{e^{0{,}36}-1} \approx 0{,}63
\]

Se comparan dos alternativas de personal:
\begin{itemize}
    \item \textbf{Alt. 1}: 6 agentes.
    \item \textbf{Alt. 2}: 7 agentes.
\end{itemize}

Métricas de interés en régimen estacionario: $W_q$ (tiempo en cola, min) y $P(W_q>2\text{ min})$ (nivel de servicio).
\end{ejemplobox}

Este sistema \textbf{no tiene solución analítica exacta}: es una cola $M/G/c$ con distribución Log-normal, para la cual no existe fórmula cerrada de $W_q$ ni de $P(W_q > t)$. La aproximación de Allen--Cunneen da una referencia aproximada:

\[
W_q^{(c)} \approx \frac{C(c,\rho)}{c\mu} \cdot \frac{1+\text{CV}_S^2}{2}
\]

\subsection{Caracterización del sistema}

La carga total del sistema es:
\[
\text{Carga} = \lambda \cdot \E[S] = \frac{45}{60} \times 6{,}55 = 4{,}91 \text{ agentes-equivalentes}
\]

Utilización teórica: $\rho_1 = 4{,}91/6 = 0{,}82$; $\rho_2 = 4{,}91/7 = 0{,}70$.

\subsection{Warm-up (Welch)}

Se ejecutan 10 réplicas largas de $T=2000$ h cada una. La media móvil con $w=50$ h se estabiliza cerca de $t_w=500$ h. Se descarta ese período en cada réplica.

\subsection{Diseño del estudio}

Se usan $n=40$ réplicas de $T=1500$ h (tras descartar los primeros 500 h de warm-up). CRN: el mismo flujo de llegadas y el mismo flujo de servicio por agente se sincronizan entre Alt. 1 y Alt. 2.

\subsection{Resultados}

\begin{center}\small
\renewcommand{\arraystretch}{1.35}
\begin{tabular}{lcccc}
\toprule
\textbf{Métrica} & \textbf{Alt. 1 ($c=6$)} & \textbf{IC 95\%} & \textbf{Alt. 2 ($c=7$)} & \textbf{IC 95\%} \\
\midrule
$W_q$ (min)        & $1{,}92$ & $[1{,}78;\;2{,}06]$ & $0{,}41$ & $[0{,}38;\;0{,}44]$ \\
$P(W_q>2\text{min})$ & $0{,}283$ & $[0{,}261;\;0{,}305]$ & $0{,}039$ & $[0{,}034;\;0{,}044]$ \\
$\rho$             & $0{,}819$ & $[0{,}815;\;0{,}823]$ & $0{,}702$ & $[0{,}699;\;0{,}705]$ \\
\bottomrule
\end{tabular}
\end{center}

\subsection{Diferencia en $W_q$ con CRN}

Con $n=40$ réplicas pareadas:
\[
\bar\Delta_{W_q} = 1{,}51 \text{ min}, \quad S_\Delta = 0{,}121 \text{ min}
\]
\[
IC_{95\%}(\Delta): \quad 1{,}51 \pm 2{,}023 \times \frac{0{,}121}{\sqrt{40}} = 1{,}51 \pm 0{,}039 = [1{,}47;\; 1{,}55] \text{ min}
\]

El IC no contiene 0; la diferencia es altamente significativa.

\subsection{Eficiencia de CRN}

\begin{center}\small
\renewcommand{\arraystretch}{1.25}
\begin{tabular}{lcc}
\toprule
\textbf{Método} & $S_\Delta$ (min) & Semiancho IC (min) \\
\midrule
Muestras independientes & $0{,}198$ & $0{,}063$ \\
CRN                     & $0{,}121$ & $0{,}039$ \\
\bottomrule
\end{tabular}
\end{center}

CRN redujo el semiancho del IC en un 38\%. Para lograr el mismo semiancho con muestras independientes se necesitarían $\lceil 40 \times (0{,}198/0{,}121)^2 \rceil = 107$ réplicas.

\subsection{Análisis de sensibilidad: ¿qué si $\lambda$ aumenta un 10\%?}

Con $\lambda' = 49{,}5$ llamadas/h, la carga sube a $5{,}40$ agentes-equivalentes. Para Alt. 1 ($c=6$): $\rho'_1 = 0{,}90$. Para Alt. 2 ($c=7$): $\rho'_2 = 0{,}77$.

Se ejecutan 20 réplicas adicionales con $\lambda' = 49{,}5$:

\begin{center}\small
\renewcommand{\arraystretch}{1.3}
\begin{tabular}{lcccc}
\toprule
 & \multicolumn{2}{c}{\textbf{Alt. 1 ($c=6$)}} & \multicolumn{2}{c}{\textbf{Alt. 2 ($c=7$)}} \\
\cmidrule(lr){2-3} \cmidrule(lr){4-5}
\textbf{Escenario} & $\bar W_q$ & IC 95\% & $\bar W_q$ & IC 95\% \\
\midrule
$\lambda = 45$ (base) & $1{,}92$ & $[1{,}78;\;2{,}06]$ & $0{,}41$ & $[0{,}38;\;0{,}44]$ \\
$\lambda = 49{,}5$ ($+10\%$) & $4{,}28$ & $[3{,}91;\;4{,}65]$ & $0{,}72$ & $[0{,}66;\;0{,}78]$ \\
\bottomrule
\end{tabular}
\end{center}

Un aumento del 10\% en $\lambda$ más que \textbf{duplica} $W_q$ para Alt. 1 ($1{,}92 \to 4{,}28$), pero solo aumenta 75\% para Alt. 2 ($0{,}41 \to 0{,}72$). Alt. 2 es mucho más robusta ante incrementos de demanda, gracias al colchón de capacidad adicional ($\rho = 0{,}70$ vs.\ $0{,}82$).

\begin{resultadobox}[Recomendación al gerente]
Con 6 agentes, casi el 28\% de las llamadas esperan más de 2 minutos, lo que incumple un estándar típico de nivel de servicio ($P(W_q>2\text{ min})\le 0{,}05$). Con 7 agentes, solo el 4\% supera ese umbral.

Adicionalmente, la Alt. 2 es significativamente más robusta ante aumentos de demanda: un incremento del 10\% en la tasa de llegadas duplicaría la espera con 6 agentes, pero causaría un aumento manejable con 7.

La simulación, con 40 réplicas y CRN, cuantifica estas diferencias con alta precisión y proporciona los intervalos de confianza necesarios para una decisión fundamentada.
\end{resultadobox}

\begin{notabox}[Resumen metodológico del Módulo 3]
El proceso completo de un estudio de simulación integra los cuatro temas:
\begin{enumerate}
    \item \textbf{T3.1 — Entradas}: ajustar distribuciones a datos (aquí: Log-normal para $S$).
    \item \textbf{T3.2 — V\&V}: verificar el código y validar contra benchmarks analíticos ($M/G/1$, Allen--Cunneen).
    \item \textbf{T3.3 — Terminal}: si el estudio fuera sobre un turno específico (horizonte finito).
    \item \textbf{T3.4 — Estacionario}: para el comportamiento de largo plazo; batch means o réplicas con warm-up; CRN para comparaciones eficientes.
\end{enumerate}
Estos pasos no son opcionales: omitir alguno compromete la credibilidad de las conclusiones.
\end{notabox}

% ══════════════════════════════════════════════════
\section*{Referencias}
% ══════════════════════════════════════════════════

\begin{enumerate}[label={[\arabic*]}]
    \item Law, A.\,M. (2015). \textit{Simulation Modeling and Analysis}, 5th ed. McGraw-Hill. Cap.~9.
    \item Banks, J., Carson, J.\,S., Nelson, B.\,L. \& Nicol, D.\,M. (2014). \textit{Discrete-Event System Simulation}, 5th ed. Pearson. Cap.~11--12.
    \item Welch, P.\,D. (1983). The statistical analysis of simulation results. En \textit{The Computer Performance Modeling Handbook}. Academic Press.
    \item Nelson, B.\,L. (2013). \textit{Foundations and Methods of Stochastic Simulation}. Springer. Cap.~9--10.
    \item Alexopoulos, C. \& Goldsman, D. (2004). To batch or not to batch? \textit{ACM Transactions on Modeling and Simulation}, 14(1), 76--114.
    \item Kim, S.\,H. \& Nelson, B.\,L. (2006). Selecting the best system. En \textit{Handbooks in Operations Research and Management Science}, Vol.~13. Elsevier.
\end{enumerate}

\end{document}
